\documentclass[12pt, a4paper]{article}
\usepackage[alf]{abntex2cite}
\usepackage[utf8]{inputenc}
\usepackage[left=1cm, top=0.1cm, bottom=0.1cm, right=1cm]{geometry}
\usepackage{graphicx}
\usepackage{setspace}
\usepackage{enumitem}
\usepackage{longtable}
\usepackage{array}
\usepackage{booktabs}
\usepackage{indentfirst}
\usepackage[portuguese]{babel}
\usepackage{caption}
\usepackage{titlesec}
\usepackage{titling}
\setlength{\droptitle}{-3cm}
\bibliographystyle{abntex2-alf}
\singlespacing

\titleformat{\section}{\normalsize\bfseries}{\thesection}{1em}{}
\titleformat{\subsection}{\small\bfseries}{\thesubsection}{1em}{}
\titlespacing*{\section}{0pt}{6pt}{3pt}
\titlespacing*{\subsection}{0pt}{4pt}{2pt}

\title{\large Estudo Dirigido 3}
\author{\small Lucas Silva de Oliveira}
\date{\small\today}

\begin{document}
\maketitle
\vspace{-0.99cm}  
\begin{center}
\small
\textbf{Instituição:} UEFS \quad | \quad
\textbf{Disciplina:} Aprendizado de Máquina e Ciência de Dados \quad | \quad
\textbf{Professor:} Dr. Rodrigo Tripodi Calumby
\end{center}

Na sua linguagem de programação e bibliotecas de apoio, investigue quais ferramentas de pré-processamento de dados estão disponíveis. 
Crie um resumo de componentes de código que permitam executar as tarefas de pré-processamento discutidas em sala: Limpeza, Integração, Redução e Transformação.

Em Python:
\begin{itemize}
    \item Pandas: É a biblioteca principal para manipulação de datasets, ela auxilia nas tarefas básicas de pré-processamento como: deletar colunas, imputar valores ausentes, filtrar dados entre outras funcionalidades.

    \item Numpy: Biblioteca que permite operações numéricas de limpeza de dados e é 
    frequentemente usado ao lado do Pandas para pré-processamento eficiente em 
    termos de desempenho.

    \item Scikit-learn: fornece um módulo de pré-processamento dedicado que 
    é muito utilizado em fluxos de trabalho de aprendizado de máquina.
\end{itemize}
\bibliography{referencias}

\end{document}  

